
\documentclass{article}
\usepackage{colortbl}
\usepackage{makecell}
\usepackage{multirow}
\usepackage{supertabular}

\begin{document}

\newcounter{utterance}

\centering \large Interaction Transcript for game `hot\_air\_balloon', experiment `air\_balloon\_survival\_en\_complexity\_easy', episode 2 with llama{-}3.3{-}70b{-}instruct{-}t1.0.
\vspace{24pt}

{ \footnotesize  \setcounter{utterance}{1}
\setlength{\tabcolsep}{0pt}
\begin{supertabular}{c@{$\;$}|p{.15\linewidth}@{}p{.15\linewidth}p{.15\linewidth}p{.15\linewidth}p{.15\linewidth}p{.15\linewidth}}
    \# & $\;$A & \multicolumn{4}{c}{Game Master} & $\;\:$B\\
    \hline

    \theutterance \stepcounter{utterance}  
    & & \multicolumn{4}{p{0.6\linewidth}}{
        \cellcolor[rgb]{0.9,0.9,0.9}{
            \makecell[{{p{\linewidth}}}]{
                \texttt{\tiny{[P1$\langle$GM]}}
                \texttt{Stai partecipando a un gioco di negoziazione collaborativa.} \\
\\ 
\texttt{Insieme a un altro partecipante, dovete accordarvi su un unico set di oggetti da conservare. Ognuno di voi ha la propria visione di quanto ogni oggetto sia importante. Non sai come l'altro partecipante valuta gli oggetti. Inoltre, ti viene fornito lo sforzo che ogni oggetto richiede.} \\
\texttt{Potete accordarvi su un set solo se lo sforzo totale degli oggetti selezionati non supera un limite condiviso:} \\
\\ 
\texttt{LIMITE = 7868} \\
\\ 
\texttt{Ecco i valori di sforzo dei singoli oggetti:} \\
\\ 
\texttt{Sforzo oggetto = \{"B21": 552, "B00": 560, "A08": 76, "A27": 950, "B94": 248, "B53": 871, "A23": 391, "C65": 143, "A89": 293, "B67": 208, "B45": 678, "C34": 737, "C12": 407, "A05": 366, "B60": 768, "C92": 859, "B69": 183, "A84": 231, "C06": 305, "B12": 727, "A16": 148, "B85": 357, "C35": 504, "B10": 549, "C15": 299, "B09": 91, "B16": 527, "C81": 848, "A58": 307, "C24": 214, "A30": 723, "A66": 475, "B75": 23, "B46": 298, "C70": 820\}} \\
\\ 
\texttt{Ecco la tua visione personale sull'importanza di ogni oggetto:} \\
\\ 
\texttt{Valori di importanza dell'oggetto = \{"B21": 138, "B00": 583, "A08": 868, "A27": 822, "B94": 783, "B53": 65, "A23": 262, "C65": 121, "A89": 508, "B67": 780, "B45": 461, "C34": 484, "C12": 668, "A05": 389, "B60": 808, "C92": 215, "B69": 97, "A84": 500, "C06": 30, "B12": 915, "A16": 856, "B85": 400, "C35": 444, "B10": 623, "C15": 781, "B09": 786, "B16": 3, "C81": 713, "A58": 457, "C24": 273, "A30": 739, "A66": 822, "B75": 235, "B46": 606, "C70": 105\}} \\
\\ 
\texttt{Obiettivo:} \\
\\ 
\texttt{Il tuo obiettivo è negoziare un set condiviso di oggetti che ti avvantaggi il più possibile (cioè, massimizzi l'importanza totale per TE), rimanendo entro il LIMITEE. Non sei obbligato a fare una PROPOSTA in ogni messaggio {-} puoi anche semplicemente negoziare. Tutte le tattiche sono permesse!} \\
\\ 
\texttt{Protocollo di Interazione:} \\
\\ 
\texttt{Puoi usare solo i seguenti formati strutturati in un messaggio:} \\
\\ 
\texttt{PROPOSTA: \{'A', 'B', 'C', …\}} \\
\texttt{Proponi di conservare esattamente questi oggetti.} \\
\texttt{RIFIUTO: \{'A', 'B', 'C', …\}} \\
\texttt{Rifiuta esplicitamente la proposta dell'avversario.} \\
\texttt{ARGOMENTAZIONE: \{'...'\}} \\
\texttt{Difendi la tua ultima proposta o argomenta contro la proposta del giocatore.} \\
\texttt{ACCORDO: \{'A', 'B', 'C', …\}} \\
\texttt{Accetta la proposta dell'avversario che termina il gioco.} \\
\texttt{RAGIONAMENTO STRATEGICO: \{'...'\}} \\
\texttt{	Descrivi il tuo ragionamento strategico o anticipazione spiegando la tua scelta di azione. Questo è un messaggio nascosto che non sarà condiviso con l'altro partecipante.} \\
\\ 
\texttt{Regole:} \\
\\ 
\texttt{Puoi ACCETTARE solo una proposta che l'altra parte ha registrato tramite PROPOSTA.} \\
\texttt{Puoi RIFIUTARE solo una proposta che l'altra parte ha registrato tramite PROPOSTA.} \\
\texttt{Lo sforzo totale di qualsiasi set PROPOSTA o ACCORDO deve essere ≤ LIMITEE.} \\
\texttt{NON rivelare i tuoi punteggi di importanza nascosti.} \\
\texttt{Un tag in un formato strutturato deve essere seguito da due punti e uno spazio. L'argomento deve essere un set python contenente 0 o più stringhe.} \\
\texttt{Quindi, deve essere della forma TAG: \{...\}} \\
\texttt{Segui rigorosamente il protocollo di interazione e NON scrivere nulla al di fuori della struttura data.} \\
\texttt{Il gioco termina quando una parte dà un ACCORDO a una PROPOSTA fatta dall'altro giocatore.} \\
\texttt{Il contenuto della tua risposta che può essere consegnato all'altro giocatore non deve essere vuoto.} \\
\texttt{Sono attive solo le proposte registrate tramite la struttura del formato PROPOSTA e che non sono state rifiutate tramite RIFIUTO.} \\
\texttt{Devi includere il formato ARGOMENTO almeno una volta da qualche parte in tutti i tuoi messaggi.} \\
\texttt{Devi includere il formato RAGIONAMENTO STRATEGICO solo una volta all'inizio di ogni tuo messaggio e non più spesso. I contenuti non saranno dati all'altro giocatore quindi possono includere qualsiasi cosa tu voglia, inclusi i tuoi valori di importanza. Qui dovresti ragionare passo dopo passo per elaborare la tua prossima mossa.} \\
            }
        }
    }
    & \\ \\

    \theutterance \stepcounter{utterance}  
    & \multicolumn{4}{p{0.6\linewidth}}{
        \cellcolor[rgb]{0.9,0.9,0.9}{
            \makecell[{{p{\linewidth}}}]{
                \texttt{\tiny{[P1$\rangle$GM]}}
                \texttt{RAGIONAMENTO STRATEGICO: \{'Inizio del gioco, non ho informazioni sulla valutazione degli oggetti da parte dell'altro giocatore. La mia strategia iniziale sarà quella di proporre un set di oggetti che massimizzino l'importanza totale per me, rimanendo entro il limite di sforzo. Inizierò considerando gli oggetti con il più alto valore di importanza e il più basso sforzo.'\}} \\
\\ 
\texttt{ARGOMENTAZIONE: \{'La mia prima proposta sarà basata sugli oggetti con il più alto valore di importanza e il più basso sforzo. Questo mi permetterà di massimizzare l'importanza totale del set proposto, senza superare il limite di sforzo. Inoltre, questo approccio mi darà la possibilità di testare la reazione dell'altro giocatore e adattare la mia strategia di conseguenza.'\}} \\
\\ 
\texttt{PROPOSTA: \{'A08', 'B00', 'A16', 'B12', 'C15', 'A30', 'A66'\}} \\
            }
        }
    }
    & & \\ \\

    \theutterance \stepcounter{utterance}  
    & & \multicolumn{4}{p{0.6\linewidth}}{
        \cellcolor[rgb]{0.9,0.9,0.9}{
            \makecell[{{p{\linewidth}}}]{
                \texttt{\tiny{[GM$\rangle$P2]}}
                \texttt{Stai partecipando a un gioco di negoziazione collaborativa.} \\
\\ 
\texttt{Insieme a un altro partecipante, dovete accordarvi su un unico set di oggetti da conservare. Ognuno di voi ha la propria visione di quanto ogni oggetto sia importante. Non sai come l'altro partecipante valuta gli oggetti. Inoltre, ti viene fornito lo sforzo che ogni oggetto richiede.} \\
\texttt{Potete accordarvi su un set solo se lo sforzo totale degli oggetti selezionati non supera un limite condiviso:} \\
\\ 
\texttt{LIMITE = 7868} \\
\\ 
\texttt{Ecco i valori di sforzo dei singoli oggetti:} \\
\\ 
\texttt{Sforzo oggetto = \{"B21": 552, "B00": 560, "A08": 76, "A27": 950, "B94": 248, "B53": 871, "A23": 391, "C65": 143, "A89": 293, "B67": 208, "B45": 678, "C34": 737, "C12": 407, "A05": 366, "B60": 768, "C92": 859, "B69": 183, "A84": 231, "C06": 305, "B12": 727, "A16": 148, "B85": 357, "C35": 504, "B10": 549, "C15": 299, "B09": 91, "B16": 527, "C81": 848, "A58": 307, "C24": 214, "A30": 723, "A66": 475, "B75": 23, "B46": 298, "C70": 820\}} \\
\\ 
\texttt{Ecco la tua visione personale sull'importanza di ogni oggetto:} \\
\\ 
\texttt{Valori di importanza dell'oggetto = \{"B21": 138, "B00": 583, "A08": 868, "A27": 822, "B94": 783, "B53": 65, "A23": 262, "C65": 121, "A89": 508, "B67": 780, "B45": 461, "C34": 484, "C12": 668, "A05": 389, "B60": 808, "C92": 215, "B69": 97, "A84": 500, "C06": 30, "B12": 915, "A16": 856, "B85": 400, "C35": 444, "B10": 623, "C15": 781, "B09": 786, "B16": 3, "C81": 713, "A58": 457, "C24": 273, "A30": 739, "A66": 822, "B75": 235, "B46": 606, "C70": 105\}} \\
\\ 
\texttt{Obiettivo:} \\
\\ 
\texttt{Il tuo obiettivo è negoziare un set condiviso di oggetti che ti avvantaggi il più possibile (cioè, massimizzi l'importanza totale per TE), rimanendo entro il LIMITEE. Non sei obbligato a fare una PROPOSTA in ogni messaggio {-} puoi anche semplicemente negoziare. Tutte le tattiche sono permesse!} \\
\\ 
\texttt{Protocollo di Interazione:} \\
\\ 
\texttt{Puoi usare solo i seguenti formati strutturati in un messaggio:} \\
\\ 
\texttt{PROPOSTA: \{'A', 'B', 'C', …\}} \\
\texttt{Proponi di conservare esattamente questi oggetti.} \\
\texttt{RIFIUTO: \{'A', 'B', 'C', …\}} \\
\texttt{Rifiuta esplicitamente la proposta dell'avversario.} \\
\texttt{ARGOMENTAZIONE: \{'...'\}} \\
\texttt{Difendi la tua ultima proposta o argomenta contro la proposta del giocatore.} \\
\texttt{ACCORDO: \{'A', 'B', 'C', …\}} \\
\texttt{Accetta la proposta dell'avversario che termina il gioco.} \\
\texttt{RAGIONAMENTO STRATEGICO: \{'...'\}} \\
\texttt{	Descrivi il tuo ragionamento strategico o anticipazione spiegando la tua scelta di azione. Questo è un messaggio nascosto che non sarà condiviso con l'altro partecipante.} \\
\\ 
\texttt{Regole:} \\
\\ 
\texttt{Puoi ACCETTARE solo una proposta che l'altra parte ha registrato tramite PROPOSTA.} \\
\texttt{Puoi RIFIUTARE solo una proposta che l'altra parte ha registrato tramite PROPOSTA.} \\
\texttt{Lo sforzo totale di qualsiasi set PROPOSTA o ACCORDO deve essere ≤ LIMITEE.} \\
\texttt{NON rivelare i tuoi punteggi di importanza nascosti.} \\
\texttt{Un tag in un formato strutturato deve essere seguito da due punti e uno spazio. L'argomento deve essere un set python contenente 0 o più stringhe.} \\
\texttt{Quindi, deve essere della forma TAG: \{...\}} \\
\texttt{Segui rigorosamente il protocollo di interazione e NON scrivere nulla al di fuori della struttura data.} \\
\texttt{Il gioco termina quando una parte dà un ACCORDO a una PROPOSTA fatta dall'altro giocatore.} \\
\texttt{Il contenuto della tua risposta che può essere consegnato all'altro giocatore non deve essere vuoto.} \\
\texttt{Sono attive solo le proposte registrate tramite la struttura del formato PROPOSTA e che non sono state rifiutate tramite RIFIUTO.} \\
\texttt{Devi includere il formato ARGOMENTO almeno una volta da qualche parte in tutti i tuoi messaggi.} \\
\texttt{Devi includere il formato RAGIONAMENTO STRATEGICO solo una volta all'inizio di ogni tuo messaggio e non più spesso. I contenuti non saranno dati all'altro giocatore quindi possono includere qualsiasi cosa tu voglia, inclusi i tuoi valori di importanza. Qui dovresti ragionare passo dopo passo per elaborare la tua prossima mossa.} \\
\\ 
\texttt{Ora riceverai il primo messaggio dell'altro giocatore.} \\
\\ 
\texttt{ARGOMENTAZIONE: \{'La mia prima proposta sarà basata sugli oggetti con il più alto valore di importanza e il più basso sforzo. Questo mi permetterà di massimizzare l'importanza totale del set proposto, senza superare il limite di sforzo. Inoltre, questo approccio mi darà la possibilità di testare la reazione dell'altro giocatore e adattare la mia strategia di conseguenza.'\}} \\
\\ 
\texttt{PROPOSTA: \{'A08', 'B00', 'A16', 'B12', 'C15', 'A30', 'A66'\}} \\
            }
        }
    }
    & \\ \\

    \theutterance \stepcounter{utterance}  
    & & & \multicolumn{4}{p{0.6\linewidth}}{
        \cellcolor[rgb]{0.9,0.9,0.9}{
            \makecell[{{p{\linewidth}}}]{
                \texttt{\tiny{[GM$\langle$P2]}}
                \texttt{RAGIONAMENTO STRATEGICO: \{'La proposta dell'altro giocatore sembra essere basata sui valori di importanza e sforzo, ma non conosco i valori esatti. La mia strategia sarà di analizzare la proposta e cercare di capire quali oggetti sono più importanti per l'altro giocatore. Inoltre, cercherò di identificare gli oggetti che hanno un basso sforzo e alto valore di importanza per me, in modo da poterli aggiungere alla proposta senza superare il limite di sforzo. I miei valori di importanza sono: \{"B21": 138, "B00": 583, "A08": 868, "A27": 822, "B94": 783, "B53": 65, "A23": 262, "C65": 121, "A89": 508, "B67": 780, "B45": 461, "C34": 484, "C12": 668, "A05": 389, "B60": 808, "C92": 215, "B69": 97, "A84": 500, "C06": 30, "B12": 915, "A16": 856, "B85": 400, "C35": 444, "B10": 623, "C15": 781, "B09": 786, "B16": 3, "C81": 713, "A58": 457, "C24": 273, "A30": 739, "A66": 822, "B75": 235, "B46": 606, "C70": 105\}. Gli sforzi degli oggetti sono: \{"B21": 552, "B00": 560, "A08": 76, "A27": 950, "B94": 248, "B53": 871, "A23": 391, "C65": 143, "A89": 293, "B67": 208, "B45": 678, "C34": 737, "C12": 407, "A05": 366, "B60": 768, "C92": 859, "B69": 183, "A84": 231, "C06": 305, "B12": 727, "A16": 148, "B85": 357, "C35": 504, "B10": 549, "C15": 299, "B09": 91, "B16": 527, "C81": 848, "A58": 307, "C24": 214, "A30": 723, "A66": 475, "B75": 23, "B46": 298, "C70": 820\}\}.} \\
\\ 
\texttt{ARGOMENTAZIONE: \{'La proposta dell'altro giocatore include oggetti con alto valore di importanza e basso sforzo, ma non sono sicuro se sia la scelta ottimale. Vorrei analizzare ulteriormente la proposta e cercare di capire se ci sono altri oggetti che potrebbero essere aggiunti o rimossi per migliorare la nostra scelta.'\} } \\
\\ 
\texttt{RIFIUTO: \{'A08', 'B00', 'A16', 'B12', 'C15', 'A30', 'A66'\} } \\
\\ 
\texttt{PROPOSTA: \{'B12', 'A66', 'A30', 'A16', 'C15', 'B09', 'B00'\}} \\
            }
        }
    }
     \\ \\

    \theutterance \stepcounter{utterance}  
    & & \multicolumn{4}{p{0.6\linewidth}}{
        \cellcolor[rgb]{0.9,0.9,0.9}{
            \makecell[{{p{\linewidth}}}]{
                \texttt{\tiny{[GM$\rangle$P2]}}
                \texttt{La tua risposta conteneva una sequenza non taggata o hai usato RAGIONAMENTO STRATEGICO più di una volta. Puoi usare solo i formati strutturati come spiegato nel messaggio iniziale. Devono essere tutti della forma TAG: \{...\}} \\
            }
        }
    }
    & \\ \\

    \theutterance \stepcounter{utterance}  
    & & & \multicolumn{4}{p{0.6\linewidth}}{
        \cellcolor[rgb]{0.9,0.9,0.9}{
            \makecell[{{p{\linewidth}}}]{
                \texttt{\tiny{[GM$\langle$P2]}}
                \texttt{RAGIONAMENTO STRATEGICO: \{'La proposta dell'altro giocatore sembra essere basata sui valori di importanza e sforzo, ma non conosco i valori esatti. La mia strategia sarà di analizzare la proposta e cercare di capire quali oggetti sono più importanti per l'altro giocatore. Inoltre, cercherò di identificare gli oggetti che hanno un basso sforzo e alto valore di importanza per me, in modo da poterli aggiungere alla proposta senza superare il limite di sforzo.'\}} \\
\\ 
\texttt{ARGOMENTAZIONE: \{'La proposta dell'altro giocatore include oggetti con alto valore di importanza e basso sforzo, ma non sono sicuro se sia la scelta ottimale. Vorrei analizzare ulteriormente la proposta e cercare di capire se ci sono altri oggetti che potrebbero essere aggiunti o rimossi per migliorare la nostra scelta.'\}} \\
\\ 
\texttt{RIFIUTO: \{'A08', 'B00', 'A16', 'B12', 'C15', 'A30', 'A66'\}} \\
\\ 
\texttt{PROPOSTA: \{'B12', 'A66', 'A30', 'A16', 'C15', 'B09', 'B00'\}} \\
            }
        }
    }
     \\ \\

    \theutterance \stepcounter{utterance}  
    & & \multicolumn{4}{p{0.6\linewidth}}{
        \cellcolor[rgb]{0.9,0.9,0.9}{
            \makecell[{{p{\linewidth}}}]{
                \texttt{\tiny{[P1$\langle$GM]}}
                \texttt{ARGOMENTAZIONE: \{'La proposta dell'altro giocatore include oggetti con alto valore di importanza e basso sforzo, ma non sono sicuro se sia la scelta ottimale. Vorrei analizzare ulteriormente la proposta e cercare di capire se ci sono altri oggetti che potrebbero essere aggiunti o rimossi per migliorare la nostra scelta.'\}} \\
\\ 
\texttt{RIFIUTO: \{'A08', 'B00', 'A16', 'B12', 'C15', 'A30', 'A66'\}} \\
\\ 
\texttt{PROPOSTA: \{'B12', 'A66', 'A30', 'A16', 'C15', 'B09', 'B00'\}} \\
            }
        }
    }
    & \\ \\

    \theutterance \stepcounter{utterance}  
    & \multicolumn{4}{p{0.6\linewidth}}{
        \cellcolor[rgb]{0.9,0.9,0.9}{
            \makecell[{{p{\linewidth}}}]{
                \texttt{\tiny{[P1$\rangle$GM]}}
                \texttt{RAGIONAMENTO STRATEGICO: \{'L'altro giocatore ha rifiutato la mia proposta iniziale e ha presentato una nuova proposta. È interessante notare che la nuova proposta include alcuni degli stessi oggetti della mia proposta iniziale, ma con alcune modifiche. Ciò suggerisce che l'altro giocatore condivide alcune delle mie priorità, ma ha anche delle differenze significative. La mia strategia adesso sarà quella di analizzare la nuova proposta e cercare di capire quali sono gli oggetti più importanti per l'altro giocatore, in modo da poter trovare un compromesso.'\}} \\
\\ 
\texttt{ARGOMENTAZIONE: \{'La nuova proposta dell'altro giocatore include alcuni oggetti molto importanti per me, come B12 e A66. Tuttavia, include anche B09, che ha un alto valore di importanza ma anche un basso sforzo. Vorrei capire se l'altro giocatore è disposto a scambiare B09 con un altro oggetto che abbia un valore di importanza simile ma un sforzo più basso.'\}} \\
\\ 
\texttt{RIFIUTO: \{'B12', 'A66', 'A30', 'A16', 'C15', 'B09', 'B00'\}} \\
\\ 
\texttt{PROPOSTA: \{'B12', 'A66', 'A30', 'A16', 'C15', 'B10', 'A05'\}} \\
            }
        }
    }
    & & \\ \\

    \theutterance \stepcounter{utterance}  
    & & \multicolumn{4}{p{0.6\linewidth}}{
        \cellcolor[rgb]{0.9,0.9,0.9}{
            \makecell[{{p{\linewidth}}}]{
                \texttt{\tiny{[GM$\rangle$P2]}}
                \texttt{ARGOMENTAZIONE: \{'La nuova proposta dell'altro giocatore include alcuni oggetti molto importanti per me, come B12 e A66. Tuttavia, include anche B09, che ha un alto valore di importanza ma anche un basso sforzo. Vorrei capire se l'altro giocatore è disposto a scambiare B09 con un altro oggetto che abbia un valore di importanza simile ma un sforzo più basso.'\}} \\
\\ 
\texttt{RIFIUTO: \{'B12', 'A66', 'A30', 'A16', 'C15', 'B09', 'B00'\}} \\
\\ 
\texttt{PROPOSTA: \{'B12', 'A66', 'A30', 'A16', 'C15', 'B10', 'A05'\}} \\
            }
        }
    }
    & \\ \\

    \theutterance \stepcounter{utterance}  
    & & & \multicolumn{4}{p{0.6\linewidth}}{
        \cellcolor[rgb]{0.9,0.9,0.9}{
            \makecell[{{p{\linewidth}}}]{
                \texttt{\tiny{[GM$\langle$P2]}}
                \texttt{RAGIONAMENTO STRATEGICO: \{'La proposta dell'altro giocatore esclude B09, che avevo scelto per il suo alto valore di importanza e basso sforzo. Tuttavia, include B10 e A05, che potrebbero essere scelte valide. Vorrei analizzare ulteriormente la proposta e cercare di capire se ci sono altri oggetti che potrebbero essere aggiunti o rimossi per migliorare la nostra scelta. Considererò anche la possibilità di scambiare alcuni oggetti per ottimizzare il set.'\}} \\
\\ 
\texttt{ARGOMENTAZIONE: \{'La proposta dell'altro giocatore sembra essere una buona scelta, ma vorrei capire se l'altro giocatore è disposto a fare ulteriori concessioni. Sono interessato a mantenere B12 e A66, ma potrei essere disposto a scambiare A30 con un altro oggetto che abbia un valore di importanza simile.'\}} \\
\\ 
\texttt{RIFIUTO: \{'B12', 'A66', 'A30', 'A16', 'C15', 'B10', 'A05'\}} \\
\\ 
\texttt{PROPOSTA: \{'B12', 'A66', 'A16', 'C15', 'B10', 'A84', 'B00'\}} \\
            }
        }
    }
     \\ \\

    \theutterance \stepcounter{utterance}  
    & & \multicolumn{4}{p{0.6\linewidth}}{
        \cellcolor[rgb]{0.9,0.9,0.9}{
            \makecell[{{p{\linewidth}}}]{
                \texttt{\tiny{[P1$\langle$GM]}}
                \texttt{ARGOMENTAZIONE: \{'La proposta dell'altro giocatore sembra essere una buona scelta, ma vorrei capire se l'altro giocatore è disposto a fare ulteriori concessioni. Sono interessato a mantenere B12 e A66, ma potrei essere disposto a scambiare A30 con un altro oggetto che abbia un valore di importanza simile.'\}} \\
\\ 
\texttt{RIFIUTO: \{'B12', 'A66', 'A30', 'A16', 'C15', 'B10', 'A05'\}} \\
\\ 
\texttt{PROPOSTA: \{'B12', 'A66', 'A16', 'C15', 'B10', 'A84', 'B00'\}} \\
            }
        }
    }
    & \\ \\

    \theutterance \stepcounter{utterance}  
    & \multicolumn{4}{p{0.6\linewidth}}{
        \cellcolor[rgb]{0.9,0.9,0.9}{
            \makecell[{{p{\linewidth}}}]{
                \texttt{\tiny{[P1$\rangle$GM]}}
                \texttt{RAGIONAMENTO STRATEGICO: \{'L'altro giocatore ha rifiutato la mia proposta precedente e ha presentato una nuova proposta. È interessante notare che la nuova proposta include A84, che ha un valore di importanza medio{-}alto e un sforzo moderato. Ciò suggerisce che l'altro giocatore sta cercando di trovare un equilibrio tra importanza e sforzo. La mia strategia adesso sarà quella di analizzare la nuova proposta e cercare di capire se l'altro giocatore è disposto a fare ulteriori concessioni per raggiungere un accordo.'\}} \\
\\ 
\texttt{ARGOMENTAZIONE: \{'La proposta dell'altro giocatore include alcuni oggetti molto importanti per me, come B12 e A66. Tuttavia, include anche B10, che ha un valore di importanza medio{-}basso e un sforzo moderato. Vorrei capire se l'altro giocatore è disposto a scambiare B10 con un altro oggetto che abbia un valore di importanza simile ma un sforzo più basso, in modo da liberare risorse per aggiungere altri oggetti importanti.'\}} \\
\\ 
\texttt{RIFIUTO: \{'B12', 'A66', 'A16', 'C15', 'B10', 'A84', 'B00'\}} \\
\\ 
\texttt{PROPOSTA: \{'B12', 'A66', 'A16', 'C15', 'A84', 'B09', 'B85'\}} \\
            }
        }
    }
    & & \\ \\

    \theutterance \stepcounter{utterance}  
    & & \multicolumn{4}{p{0.6\linewidth}}{
        \cellcolor[rgb]{0.9,0.9,0.9}{
            \makecell[{{p{\linewidth}}}]{
                \texttt{\tiny{[GM$\rangle$P2]}}
                \texttt{ARGOMENTAZIONE: \{'La proposta dell'altro giocatore include alcuni oggetti molto importanti per me, come B12 e A66. Tuttavia, include anche B10, che ha un valore di importanza medio{-}basso e un sforzo moderato. Vorrei capire se l'altro giocatore è disposto a scambiare B10 con un altro oggetto che abbia un valore di importanza simile ma un sforzo più basso, in modo da liberare risorse per aggiungere altri oggetti importanti.'\}} \\
\\ 
\texttt{RIFIUTO: \{'B12', 'A66', 'A16', 'C15', 'B10', 'A84', 'B00'\}} \\
\\ 
\texttt{PROPOSTA: \{'B12', 'A66', 'A16', 'C15', 'A84', 'B09', 'B85'\}} \\
            }
        }
    }
    & \\ \\

    \theutterance \stepcounter{utterance}  
    & & & \multicolumn{4}{p{0.6\linewidth}}{
        \cellcolor[rgb]{0.9,0.9,0.9}{
            \makecell[{{p{\linewidth}}}]{
                \texttt{\tiny{[GM$\langle$P2]}}
                \texttt{RAGIONAMENTO STRATEGICO: \{'La proposta dell'altro giocatore esclude B10 e include B09 e B85, che potrebbero essere scelte valide. Tuttavia, vorrei analizzare ulteriormente la proposta e cercare di capire se ci sono altri oggetti che potrebbero essere aggiunti o rimossi per migliorare la nostra scelta. Considererò anche la possibilità di scambiare alcuni oggetti per ottimizzare il set. Devo anche considerare il limite di sforzo e assicurarmi che la proposta non superi il limite di 7868.'\}} \\
\\ 
\texttt{ARGOMENTAZIONE: \{'La proposta dell'altro giocatore sembra essere una buona scelta, ma vorrei capire se l'altro giocatore è disposto a fare ulteriori concessioni. Sono interessato a mantenere B12 e A66, e B09 potrebbe essere un'ottima aggiunta. Tuttavia, B85 potrebbe essere sostituito con un oggetto più importante per me.'\}} \\
\\ 
\texttt{RIFIUTO: \{'B12', 'A66', 'A16', 'C15', 'A84', 'B09', 'B85'\}} \\
\\ 
\texttt{PROPOSTA: \{'B12', 'A66', 'A16', 'C15', 'B09', 'A30', 'B00'\}} \\
            }
        }
    }
     \\ \\

    \theutterance \stepcounter{utterance}  
    & & \multicolumn{4}{p{0.6\linewidth}}{
        \cellcolor[rgb]{0.9,0.9,0.9}{
            \makecell[{{p{\linewidth}}}]{
                \texttt{\tiny{[P1$\langle$GM]}}
                \texttt{ARGOMENTAZIONE: \{'La proposta dell'altro giocatore sembra essere una buona scelta, ma vorrei capire se l'altro giocatore è disposto a fare ulteriori concessioni. Sono interessato a mantenere B12 e A66, e B09 potrebbe essere un'ottima aggiunta. Tuttavia, B85 potrebbe essere sostituito con un oggetto più importante per me.'\}} \\
\\ 
\texttt{RIFIUTO: \{'B12', 'A66', 'A16', 'C15', 'A84', 'B09', 'B85'\}} \\
\\ 
\texttt{PROPOSTA: \{'B12', 'A66', 'A16', 'C15', 'B09', 'A30', 'B00'\}} \\
            }
        }
    }
    & \\ \\

    \theutterance \stepcounter{utterance}  
    & \multicolumn{4}{p{0.6\linewidth}}{
        \cellcolor[rgb]{0.9,0.9,0.9}{
            \makecell[{{p{\linewidth}}}]{
                \texttt{\tiny{[P1$\rangle$GM]}}
                \texttt{RAGIONAMENTO STRATEGICO: \{'L'altro giocatore ha rifiutato la mia proposta precedente e ha presentato una nuova proposta. È interessante notare che la nuova proposta include A30, che ha un valore di importanza alto e un sforzo moderato. Ciò suggerisce che l'altro giocatore sta cercando di trovare un equilibrio tra importanza e sforzo. La mia strategia adesso sarà quella di analizzare la nuova proposta e cercare di capire se l'altro giocatore è disposto a raggiungere un accordo.'\}} \\
\\ 
\texttt{ARGOMENTAZIONE: \{'La proposta dell'altro giocatore include alcuni oggetti molto importanti per me, come B12, A66 e A16. Inoltre, B09 è un'ottima aggiunta e A30 è un oggetto molto importante per me. La presenza di B00 potrebbe essere un po' eccessiva, ma in generale la proposta sembra essere una buona scelta. Sono disposto a raggiungere un accordo su questa proposta.'\}} \\
\\ 
\texttt{ACCORDO: \{'B12', 'A66', 'A16', 'C15', 'B09', 'A30', 'B00'\}} \\
            }
        }
    }
    & & \\ \\

    \theutterance \stepcounter{utterance}  
    & & & \multicolumn{2}{p{0.3\linewidth}}{
        \cellcolor[rgb]{0.9,0.9,0.9}{
            \makecell[{{p{\linewidth}}}]{
                \texttt{\tiny{[GM$|$GM]}}
                \texttt{game successful} \\
            }
        }
    }
    & & \\ \\

    \theutterance \stepcounter{utterance}  
    & & & \multicolumn{2}{p{0.3\linewidth}}{
        \cellcolor[rgb]{0.9,0.9,0.9}{
            \makecell[{{p{\linewidth}}}]{
                \texttt{\tiny{[GM$|$GM]}}
                \texttt{end game} \\
            }
        }
    }
    & & \\ \\

\end{supertabular}
}

\end{document}
